\documentclass[11pt,letterpaper]{article}
\usepackage[T1]{fontenc}
\usepackage{lmodern}
\usepackage[margin=1in,headheight=15pt]{geometry}
\usepackage{amsmath,amssymb,amsthm,mathtools}
\usepackage{microtype,booktabs,array,longtable}
\usepackage[dvipsnames]{xcolor}
\usepackage[most]{tcolorbox}
\usepackage{enumitem,fancyhdr,titlesec,listings}
\usepackage{xurl}
\usepackage{hyperref}
\definecolor{Ink}{HTML}{16324F}
\definecolor{Accent}{HTML}{167D8D}
\definecolor{Pale}{HTML}{EDF5F7}
\definecolor{Muted}{HTML}{536777}
\hypersetup{colorlinks=true,linkcolor=Ink,citecolor=Accent,urlcolor=Accent,
 pdftitle={Throwback Dynamics with Repeated Entries},
 pdfauthor={Research manuscript prepared with ChatGPT},
 pdfsubject={Recurrence classification, finite cycles, perfect mixing, and extremal periods}}
\pagestyle{fancy}\fancyhf{}
\fancyhead[L]{\small\color{Muted}Throwback dynamics with repeated entries}
\fancyhead[R]{\small\color{Muted}Research manuscript}
\fancyfoot[C]{\thepage}
\renewcommand{\headrulewidth}{0.3pt}
\titleformat{\section}{\Large\bfseries\color{Ink}}{\thesection}{0.7em}{}
\titleformat{\subsection}{\large\bfseries\color{Ink}}{\thesubsection}{0.7em}{}
\setlength{\parskip}{0.35em}
\setlength{\parindent}{1.15em}
\setlist{nosep,leftmargin=1.5em}
\newtheorem{theorem}{Theorem}[section]
\newtheorem{lemma}[theorem]{Lemma}
\newtheorem{proposition}[theorem]{Proposition}
\newtheorem{corollary}[theorem]{Corollary}
\theoremstyle{definition}
\newtheorem{definition}[theorem]{Definition}
\newtheorem{example}[theorem]{Example}
\newtheorem{remark}[theorem]{Remark}
\newcommand{\N}{\mathbb N}
\newcommand{\Z}{\mathbb Z}
\newcommand{\Q}{\mathbb Q}
\newcommand{\pos}{\operatorname{pos}}
\newcommand{\den}{\operatorname{den}}
\newcommand{\lcm}{\operatorname{lcm}}
\newcommand{\supp}{\operatorname{supp}}
\newcommand{\wt}{\operatorname{wt}}
\newcommand{\core}{\mathcal R}
\newcommand{\states}{\mathcal C}
\newcommand{\eps}{\varepsilon}
\newcommand{\ind}{\mathbf 1}
\newtcolorbox{resultbox}{colback=Pale,colframe=Accent,boxrule=0.6pt,
 arc=1.5mm,left=3mm,right=3mm,top=2mm,bottom=2mm,breakable}
\lstset{basicstyle=\small\ttfamily,breaklines=true,columns=fullflexible,
 frame=single,rulecolor=\color{Muted},backgroundcolor=\color{Pale},
 keepspaces=true,showstringspaces=false}
\allowdisplaybreaks[2]
\begin{document}
\begin{titlepage}
\thispagestyle{empty}
{\color{Accent}\sffamily\bfseries RESEARCH IN INTEGER-SEQUENCE DYNAMICS\par}
\vspace{0.7in}
{\Huge\bfseries\color{Ink}Throwback Dynamics\\[0.15em]with Repeated Entries\par}
\vspace{0.25in}
{\Large\color{Muted}Recurrence classification, exact cycles,\\perfect mixing, and extremal periods\par}
\vspace{0.3in}
{\large 19 September 2026\par}
\vspace{0.35in}
\begin{abstract}
A throwback move removes the first entry of an infinite queue and reinserts it
at the zero-based position equal to its positive integer value. We resolve
the repeated-entry recurrence question in an occurrence-labeled formulation,
and also identify exactly which numerical values recur. Every occurrence is
recurrent precisely when, for every $m$, at most $m$ entries have value at
most $m$. When this condition fails, the first overloaded initial prefix
identifies the complete recurrent core; the remaining tail never participates.
A sparse backward-parking algorithm extracts this core without simulating the
queue. We then determine every cycle length of the bounded configuration
system, obtain an explicit coprimality criterion for perfect mixing, and prove
that the largest labeled period of a recurrent core of size $r$ is
$\max_{2\le k\le r}k^{r-k+1}$. This gives a throwback interpretation of OEIS
A003320. Catalan and prime-parking-function enumerations, rational generating
functions, a Lambert-$W$ asymptotic, and an exact example of escaped limiting
density complete the analysis. All main statements have proofs; executable
standard-library Python checks accompany the manuscript.
\end{abstract}
\vspace{0.15in}
\begin{resultbox}
\textbf{Scope and status.} The starting questions are explicitly raised by
Darnell and Dribus in their 2025 paper. This manuscript presents proposed
resolutions with self-contained proofs, not a claim of established publication
priority. Literature checking did not locate a later resolution, but was not
exhaustive. The proofs have not been independently refereed or formalized in a
proof assistant. Computation is used as an additional check, not as a substitute
for proof.
\end{resultbox}
\vfill
{\small\color{Muted}Companion archive: \texttt{code/}, \texttt{data/},
\texttt{notes/}, this \LaTeX\ source, and the compiled PDF. No external Python
packages are required.}
\end{titlepage}
\pagenumbering{roman}
{\small\setlength{\parskip}{0pt}\tableofcontents}
\newpage
\pagenumbering{arabic}

\section{The problem and its provenance}\label{sec:provenance}
Darnell and Dribus prove recurrence for distinct positive inputs and leave the
repeated-entry case unresolved immediately after their Theorem 2
\cite[p.~4]{DD}. They also ask for a characterization of perfect mixing
\cite[p.~16]{DD}. These are the two questions addressed here. Their frequency
and placement results for distinct inputs provide useful context, but we give
complete proofs of the extensions used below.

The motivating OEIS entry A357081 records the leader sequence obtained from
$(3,4,5,\ldots)$ \cite{A357081}. Our target is not its original recurrence
conjecture, which is already covered by the 2025 theorem. Instead, the
central issue is the effect of repetitions: equal values can form an active
front while all later entries become inaccessible.

A necessary distinction is easy to overlook. In a queue consisting entirely
of $1$'s, the \emph{numerical value} $1$ recurs forever, but only the first two
\emph{occurrences} ever lead. We therefore retain occurrence labels throughout.
The resulting theorem determines numerical recurrence too, without confusing
these two meanings.

The principal contributions are the first-obstruction theorem
(Theorem~\ref{thm:obstruction}), its algorithmic implementation
(Theorem~\ref{thm:algorithm}), the exact cycle theorem
(Theorem~\ref{thm:cycles}), and the extremal-period theorem
(Theorem~\ref{thm:extremal}). The perfect-mixing criterion is
Corollary~\ref{cor:mixing}. The familiar sequences that appear in the enumeration
and extremal formula are not claimed to be new sequences.

\subsection{What the source check establishes}
The questions were checked against the actual 2025 article, including the
relevant PDF pages, rather than inferred from an OEIS ``conjecture'' tag.
Public searches and the relevant OEIS pages were consulted on 19 September
2026. No later treatment of the two selected questions was located. This
negative search result is a limitation of the literature check, not a proof
that no earlier or unpublished solution exists. The mathematical content
below is intended to be checked independently of that priority question.

\section{Labeled queues and closed fronts}\label{sec:model}
Let $X=(x_0,x_1,\ldots)$ be an infinite sequence with $x_i\in\Z_{>0}$.
Token $i$ has permanent weight $x_i$. A queue state is an ordering of these
tokens; its positions are numbered $0,1,2,\ldots$. If the current head is token
$i$, one move removes it and reinserts it at position $x_i$. Thus the numerical
operation is
\begin{equation}\label{eq:tau}
 \tau(x_0,x_1,\ldots)
   =(x_1,\ldots,x_{x_0},x_0,x_{x_0+1},\ldots).
\end{equation}
Let $Q_t$ denote the state before move $t$, let $L_t$ be its head token, and
let $T_t=x_{L_t}$. The sequences $(L_t)$ and $(T_t)$ are respectively the
\emph{labeled} and \emph{numerical} leader words.

A token is \emph{recurrent} when it occurs infinitely often in $(L_t)$. A
numerical value is recurrent when it occurs infinitely often in $(T_t)$.
The distinction matters whenever weights repeat.

\begin{definition}
A finite multiset $A$ of positive weights is \emph{admissible} if
\begin{equation}\label{eq:admissible}
 C_A(m):=\#\{a\in A:a\le m\}\le m\qquad(m\ge1).
\end{equation}
An infinite input is admissible when the same inequalities hold with all
its occurrences counted. Counts may initially be infinite; an admissible
input necessarily has finite multiplicity below every fixed bound.
\end{definition}
If the sorted weights of a finite multiset are
$a_0\le\cdots\le a_{n-1}$, then admissibility is equivalent to
\begin{equation}\label{eq:sorted-admissible}
 a_i\ge i+1\qquad(0\le i<n).
\end{equation}
Indeed, $a_i\le i$ exhibits at least $i+1$ weights at most $i$.
Conversely, $C_A(m)>m$ forces $a_m\le m$.
Admissibility is inherited by submultisets.

\begin{definition}
A \emph{closed front of size $r$} is a state whose first $r$ tokens all have
weight at most $r-1$. A multiset of $r\ge2$ positive weights is
\emph{critical} if
\begin{equation}\label{eq:critical}
 \max A=r-1,\qquad C_A(m)\le m\quad(1\le m\le r-2).
\end{equation}
\end{definition}
A closed front stays closed: every subsequent move is wholly inside those
$r$ positions. Tokens behind it remain at exactly the same positions, not
merely at positions with a limiting distribution.

Two elementary observations will be used repeatedly. A token not currently
leading can only stay put or move one position left. Once token $i$ has been
thrown, its position is always at most $x_i$. Also, the relative order of
never-yet-leading tokens never changes. In particular, first appearances as
leader occur in original label order $0,1,2,\ldots$, possibly stopping at a
finite label.

\section{A quantitative waiting-or-barrier lemma}\label{sec:waiting}
The following lemma is useful because it does not assume admissibility.
Its obstruction is an actual closed front, rather than an unspecified failure
of recurrence.

\begin{lemma}[Waiting or a barrier]\label{lem:waiting}
Suppose a token $A$ is currently at position $p$. Within $2^p-1$ moves, one
of the following has happened: $A$ has been a leader, or a closed front of
size at most $p$, lying strictly before $A$, has formed. In the second case
$A$ has not yet led during this waiting interval and can never lead afterward.
\end{lemma}
\begin{proof}
Induct on $p$. At $p=0$ the token is already the leader. Suppose $p>0$.
If all $p$ tokens before $A$ have weight at most $p-1$, they form the required
closed front at once. Otherwise choose a token $B$ before $A$ whose weight
is at least $p$. Its position is at most $p-1$.

Apply the induction hypothesis to $B$. A barrier before $B$ is also before
$A$, since before either token first leads their relative order cannot change.
Alternatively, $B$ leads after at most $2^{p-1}-1$ moves. At that time $A$ has
position at most $p$. Throwing $B$ shifts $A$ at least one position left,
because $B$'s weight is at least $p$. Now $A$ is at position at most $p-1$.
A second application of induction needs at most another $2^{p-1}-1$ moves.
The total is at most
\[
 (2^{p-1}-1)+1+(2^{p-1}-1)=2^p-1.
\]
Every barrier supplied by this argument excludes $A$ and is permanent.
\end{proof}

\begin{corollary}\label{cor:no-barrier}
An admissible infinite input has no closed front at any time. Every token
initially at position $j$ first leads by time $2^j-1$. A recurrent visit by a
token of weight $v$ is followed by another within at most $2^v$ moves.
\end{corollary}
\begin{proof}
A closed front of size $r$ would contain $r$ weights at most $r-1$, contrary
to admissibility. Apply Lemma~\ref{lem:waiting} initially, and then after each
throw, when the token has position $v$. The extra move for the throw changes
the bound $2^v-1$ to $2^v$ between leader times.
\end{proof}
The numerical bound is deliberately simple, not generally optimal. Its
importance here is that it gives finite-time conclusions with no probabilistic
or compactness argument.

\section{The first overloaded prefix determines everything}\label{sec:classification}
For $j\ge0$, let $A_j$ be the multiset of weights in the original prefix
$(x_0,\ldots,x_j)$. Write $A_{-1}=\varnothing$.

\begin{theorem}[First-obstruction classification]\label{thm:obstruction}
Exactly one of the following alternatives holds.

\textbf{Admissible alternative.} If
\begin{equation}\label{eq:global-condition}
 \boxed{\quad\#\{i\ge0:x_i\le m\}\le m\quad\text{for every }m\ge1,\quad}
\end{equation}
then every labeled token leads infinitely often.

\textbf{Obstructed alternative.} Otherwise let
\begin{align}
 J&=\min\{j\ge0:A_j\text{ is not admissible}\},\label{eq:J}\\
 m&=\min\{k\ge1:C_{A_J}(k)>k\}.\label{eq:m}
\end{align}
Then $x_J\le m$, $C_{A_J}(m)=m+1$, and the recurrent token set is exactly
\begin{equation}\label{eq:core}
 \boxed{\quad\core=\{i\le J:x_i\le m\}.\quad}
\end{equation}
Precisely the labels $0,1,\ldots,J$ ever lead. By the first appearance of
label $J$, and in particular by time $2^J-1$, the tokens in $\core$ occupy a
closed front of size $m+1$. Every position behind that front is then fixed.
The multiset of recurrent weights is critical. These conclusions, and the
entire leader word, are independent of the continuation after $x_J$.
\end{theorem}

\begin{proof}
If \eqref{eq:global-condition} holds, Corollary~\ref{cor:no-barrier} proves
recurrence. If it fails at a threshold $k$, finitely many occurrences already
witness $k+1$ entries of weight at most $k$. Thus a finite $J$ in
\eqref{eq:J} exists.

The multiset $A_{J-1}$ is admissible. Adding just one entry can increase any
count by only one. Therefore the least violated threshold $m$ satisfies
\[
 C_{A_J}(m)=m+1,\qquad C_{A_{J-1}}(m)=m,\qquad x_J\le m.
\]
For each $k<m$, minimality of $m$ gives $C_{A_J}(k)\le k$. Consequently
$\core$ has $m+1$ tokens, all of weight at most $m$, and its counts below
$m$ satisfy the critical inequalities. At least one weight is $m$:
otherwise $m+1$ weights would already be at most $m-1$. Thus its multiset
is critical.

We next show that every label $j\le J$ actually leads. Until its first
appearance, every token in front of $j$ has original label less than $j$:
no later label can pass an unvisited token by becoming a leader first.
A closed front strictly before $j$ would therefore be a submultiset of
$A_{j-1}$. But $A_{j-1}$ is admissible, so such a front is impossible.
Lemma~\ref{lem:waiting}, applied to the initial position $j$, proves that
$j$ first leads by time $2^j-1$.

At the first leader time of $J$, every earlier token has already led and
been thrown. The $m$ earlier tokens in $\core$ therefore lie in positions
at most $m$. Token $J$ is at position zero and also belongs to $\core$.
These $m+1$ distinct tokens fill all positions $0,\ldots,m$, so the front
is closed already at that time. No later original label has yet led, and
none can lead afterward. This proves the assertion about all first appearances
and the permanent tail.

It remains to check that every member of $\core$ is recurrent, rather than
merely trapped in a larger closed front. A proper closed subfront of size
$r<m+1$ would require $r$ core weights at most $r-1$, violating criticality.
Thus no proper barrier can occur inside the core. Repeated application of
Lemma~\ref{lem:waiting} gives recurrence of every core token.

Finally, the proof applies to every continuation of the prefix $A_J$.
No token from that continuation ever leads, so its weight is never used to
determine a move. Tracking the positions of labels $0,\ldots,J$ therefore
produces the same leader word for every continuation.
\end{proof}

\begin{corollary}[Sharp recurrence criterion]\label{cor:iff}
For an infinite positive input, the following are equivalent: every token
is recurrent; every token leads at least once; every finite initial prefix
is admissible; and \eqref{eq:global-condition} holds.
\end{corollary}
\begin{proof}
The only non-immediate implication is that every token leads at least once
only in the admissible alternative. In the obstructed alternative,
Theorem~\ref{thm:obstruction} says that no label larger than $J$ ever leads.
\end{proof}

\begin{corollary}[Numerical recurrence]\label{cor:numerical}
In the admissible alternative every initial numerical value recurs. In the
obstructed alternative a value $v$ recurs if and only if $v=x_i$ for some
$i\in\core$. Thus all initial numerical values recur exactly when the full
input's set of values is contained in the core's set of values. In particular,
for inputs with unbounded values, recurrence of every numerical value is
equivalent to \eqref{eq:global-condition}.
\end{corollary}
\begin{proof}
In the obstructed alternative only finitely many noncore leader occurrences
precede permanent capture. Each core token is recurrent. A finite core cannot
represent an unbounded set of values.
\end{proof}

\subsection{Why the first prefix, rather than the global smallest value, matters}
The recurrence criterion for \emph{all} tokens is independent of their order.
When recurrence fails, the identity of the surviving core need not be.
For example,
\[
 (2,2,2,1,1,6,7,\ldots)
 \quad\hbox{and}\quad
 (1,1,2,2,2,6,7,\ldots)
\]
have the same multiset. The first queue keeps its three initial $2$-tokens
active, while the second keeps its two initial $1$-tokens active. Looking only
for the globally least overloaded threshold would incorrectly predict the
first queue's core. The prefix in \eqref{eq:J} is essential.

\section{Examples, including a nonlocal-looking obstruction}\label{sec:examples}
For
\[
 X=(2,2,3,4,5,\ldots),
\]
one has $C_X(1)=0$ and $C_X(m)=m$ for $m\ge2$. Repetition is harmless:
every occurrence recurs. In contrast, a third initial $2$ closes the first
three positions immediately.

Consider the more instructive input
\begin{equation}\label{eq:example100}
 X=(2,100,2,1,1,\ldots).
\end{equation}
Its first three weights $(2,100,2)$ are admissible. Adding the first $1$
creates the first bad prefix, with $J=3$ and $m=2$. Hence labels $0,2,3$
are exactly the recurrent core. Label $1$, of weight $100$, appears only
transiently; label $4$, though also of weight $1$, never leads.

The initial labeled leaders are
\[
 0,1,2,0,3,2,3,0,3,2,3,0,\ldots.
\]
The numerical word is
\[
 2,100,2,2,1,2,1,2,1,2,1,2,\ldots.
\]
The eventual labeled period is $4$, while the numerical period is $2$.
This difference is not cosmetic: cycle formulas in this manuscript always
refer to labeled tokens unless explicitly stated otherwise.

For $X=(1,2,3,1,\ldots)$, the first bad prefix has $J=3$, $m=1$, and
$\core=\{0,3\}$. Thus the transiently encountered values $2$ and $3$ do not
recur. For a constant input $(1,1,1,\ldots)$ the same mechanism leaves only
labels $0,1$, although its only numerical value is recurrent.

\section{A sparse algorithm that avoids dynamical simulation}\label{sec:algorithm}
The inequalities defining admissibility are a nested matching condition.
Think of each weight $v$ as a deadline for a unit job that must occupy a
distinct positive slot not exceeding $v$.

\subsection{Backward parking}
Process tokens in original order. Assign each new token of weight $v$ the
largest currently unoccupied slot at most $v$. Fail if no such slot exists.
This is a static allocation procedure, not the throwback process itself.

\begin{lemma}\label{lem:parking}
Backward parking succeeds on a finite word if and only if its multiset is
admissible. Its first failure occurs at exactly the index $J$ in
Theorem~\ref{thm:obstruction}.
\end{lemma}
\begin{proof}
Any successful allocation gives an injection of the tokens of weight at most
$m$ into slots $1,\ldots,m$, proving necessity.

Suppose a token of weight $v$ fails while all earlier tokens were parked.
All slots $1,\ldots,v$ are occupied. Let $M+1$ be the first unoccupied slot
larger than $v$. It exists because only finitely many slots have been used.
Every slot $1,\ldots,M$ is occupied. Every token parked in those slots has
weight at most $M$: otherwise the still-free slot $M+1$ would have been
available when it parked, contradicting the choice of the largest free slot.
Together with the failed token, there are therefore $M+1$ weights at most
$M$. Failure certifies inadmissibility. The two implications also identify
the first failure.
\end{proof}

\subsection{Recovering the least threshold without sorting}
A first failure identifies $J$, but the proof above may give a threshold
larger than the least one. The following closure computation recovers the
exact $m$ in linear additional work.

Let $d(p)$ be the weight of the previously parked token occupying slot $p$.
Start with $m=v=x_J$. Scan slots $p=1,2,\ldots,m$ in order, extending the
scan whenever necessary by replacing $m$ with $\max(m,d(p))$. Stop when the
scan reaches the current $m$. The scan never reaches an empty slot: the
$M$ from Lemma~\ref{lem:parking} bounds all extensions.

At termination, the first $m$ slots are occupied by $m$ tokens of weight at
most $m$. Every earlier token of weight at most $m$ must have parked in
these slots, so $C_{A_{J-1}}(m)=m$. Conversely, if a threshold $k\ge v$
satisfies that equality, its first $k$ slots are occupied precisely by weights
at most $k$, and the closure cannot grow past $k$. Hence the terminal value
is the least saturated threshold at least $v$, exactly the least threshold
violated by adding $v$. The core is the set of occupants of slots
$1,\ldots,m$, together with the failed token.

\begin{theorem}[Static extraction and certification]\label{thm:algorithm}
Given a prefix through its first obstruction $J$, backward parking and the
closure scan compute $J$, $m$, and the complete recurrent core. Using a
union-by-size disjoint-set structure on occupied intervals, this takes
$O((J+1)\alpha(J+1))$ amortized data-structure operations and $O(J+1)$ space.
The bound uses unit-cost integer and dictionary operations; it is not a
bit-complexity bound for arbitrarily large weights.
\end{theorem}
\begin{proof}
The preceding lemmas prove correctness. The occupied slots form disjoint
intervals. If $v$ is free, it is the desired slot; otherwise let $[a,b]$ be
the occupied component containing $v$. The desired slot is $a-1$, with
$a=1$ meaning failure. Occupying one slot joins at most two neighboring
components. Union by size, path compression, and a stored component minimum
support these operations in the stated amortized bound. The final closure
visits at most $J$ occupied slots.
\end{proof}

The distributed implementation uses sparse dictionaries keyed by occupied
slots, so a weight such as $10^{40}$ does not allocate $10^{40}$ positions.
The certificate has a deliberately independent verifier: it sorts the prefix
before $J$ and checks \eqref{eq:sorted-admissible}, then checks the least
failure in the sorted extended prefix and the selected labels. A user need
not trust the parking implementation to verify its output.

\begin{lstlisting}[caption={High-level static classification algorithm.}]
for each token j with weight v:
    p = largest free positive slot <= v
    if p exists:
        park token j in p
    else:
        m = v
        cursor = 1
        while cursor <= m:
            m = max(m, weight of occupant of slot cursor)
            cursor += 1
        return (J=j, threshold=m,
                core={occupants of slots 1,...,m} union {j})
\end{lstlisting}
On a finite admissible input, ``no obstruction found'' certifies only that
inspected prefix. On an infinite admissible stream the search does not halt.
This is an unavoidable distinction between finite evidence and an assertion
about an arbitrary infinite continuation.

\section{Counting all possible recurrent cores}\label{sec:enumeration}
A critical multiset is not merely necessary: any ordering of it, followed
by any infinite tail, has that multiset as its complete recurrent core.
Indeed, the initial front is closed and no proper closed subfront is possible.
Thus critical multisets are the exact finite possibilities.

\subsection{A cyclic-minimum lemma}
\begin{lemma}\label{lem:cycle}
Let $a_1,\ldots,a_N$ be integers with sum $1$. Exactly one of their $N$
cyclic starting positions gives strictly positive partial sums.
\end{lemma}
\begin{proof}
Put $S_0=0$ and $S_j=a_1+\cdots+a_j$. Start just after the last occurrence
of the minimum of $S_0,\ldots,S_{N-1}$. Unwrapped partial sums are positive
by this last-minimum choice. Wrapped partial sums have the form
$1+S_j-S_k$ and are positive because the $S_j$ are integers and
$S_j\ge S_k$. Conversely, positivity of an unwrapped segment requires all
later $S_j$ to exceed $S_k$, while positivity after wrapping requires all
earlier $S_j$ to be at least $S_k$. Thus any successful start is precisely
that last minimum.
\end{proof}

\begin{theorem}[Core enumeration]\label{thm:counts}
For $r\ge2$, the number of critical ordered weight vectors of length $r$ is
\begin{equation}\label{eq:orderedcount}
 E_r=(r-1)^{r-1},
\end{equation}
and the number of critical multisets is
\begin{equation}\label{eq:multisetcount}
 H_r=\frac1r\binom{2r-2}{r-1}=C_{r-1}.
\end{equation}
Consequently a uniformly chosen closed weight vector from
$\{1,\ldots,r-1\}^r$ has all $r$ tokens recurrent with probability
$1/(r-1)$.
\end{theorem}
\begin{proof}
Set $N=r-1$ and replace each weight $x$ by $b=r-x\in\{1,\ldots,N\}$.
If $c_k$ is the number of these $b$'s equal to $k$, criticality becomes
\begin{equation}\label{eq:prime-parking}
 c_1+\cdots+c_k>k\qquad(1\le k\le N),
 \qquad c_1+\cdots+c_N=N+1.
\end{equation}
For example, the count in \eqref{eq:prime-parking} is
$r-C_A(r-k-1)$, so the equivalence follows directly from
\eqref{eq:critical}.

Apply Lemma~\ref{lem:cycle} to $c_k-1$, whose sum is $1$. Adding a common
residue modulo $N$ to all $b$-values cyclically rotates the counts. Every
orbit of ordered words has size $N$, and exactly one word in it satisfies
\eqref{eq:prime-parking}. There are $N^{N+1}$ words, giving
$N^{N+1}/N=N^N$ critical ordered vectors.

For multisets, count weak compositions of $N+1$ into $N$ parts. There are
$\binom{2N}{N-1}$ such vectors. A count vector cannot have a nontrivial
cyclic period: the number of repetitions would divide both $N$ and $N+1$.
Thus all its cyclic orbits have size $N$, again with exactly one valid
rotation. Their number is
\[
 \frac1N\binom{2N}{N-1}
   =\frac1{N+1}\binom{2N}{N}.
\]
The probability statement follows by dividing \eqref{eq:orderedcount} by
$(r-1)^r$.
\end{proof}
The transformed objects in \eqref{eq:prime-parking} are prime parking
functions. Their known count is recorded in A000312; $C_n$ is A000108
\cite{A000312,A000108}. The contribution here is their identification with
throwback recurrent cores, not a new enumeration of parking functions or
Catalan objects.

\subsection{Generating functions and asymptotics}
Let $C(z)=(1-\sqrt{1-4z})/(2z)$. Formula~\eqref{eq:multisetcount} gives
\begin{equation}\label{eq:core-ogf}
 \sum_{r\ge2}H_rz^r=z(C(z)-1)
  =\frac{1-\sqrt{1-4z}}2-z.
\end{equation}
In particular, $C_0=1$ and
$C_n=\sum_{j=0}^{n-1}C_jC_{n-1-j}$ give a recurrence for $H_r$.
Stirling's formula applied to \eqref{eq:multisetcount} yields
\[
 H_r\sim\frac{4^{r-1}}{\sqrt\pi\,(r-1)^{3/2}}.
\]

For the ordered count, let $U(z)$ be the formal solution of $U=ze^U$;
equivalently, $U=-W(-z)$ on the branch analytic at zero. Then
\begin{equation}\label{eq:core-egf}
 \sum_{r\ge2} E_r\frac{z^r}{r!}=1-e^{-U(z)}-z.
\end{equation}
For completeness, Lagrange coefficient extraction gives
\[
 [z^r](1-e^{-U(z)})
  =\frac1r[u^{r-1}]e^{-u}e^{ru}
  =\frac{(r-1)^{r-1}}{r!}.
\]
The $r=1$ term is $z$, which is removed because a positive-weight core has
at least two tokens. Thus the formal $0^0$ convention used for that auxiliary
coefficient does not introduce a size-one positive core.

\section{Finite bounded configurations and exact cycles}\label{sec:cycles}
The recurrence classification tells us which tokens survive. We now determine
how the surviving or selected tokens move.

\subsection{The autonomous state space}
Select $r$ labeled tokens whose weights, in nondecreasing order, are
$s_0,\ldots,s_{r-1}$. Assume either that this multiset is admissible and every
unselected token has weight at least $s_{r-1}$, or that it is a critical closed
core. Set
\begin{equation}\label{eq:q}
 q_i=s_i-i+1.
\end{equation}
In the admissible case all $q_i\ge2$. In the critical case all but the last
are at least $2$, and $q_{r-1}=1$.

After each selected token has been thrown, its position $p_i$ lies in
$[0,s_i]$. The bounded state space is
\begin{equation}\label{eq:states}
 \states(s)=\{(p_0,\ldots,p_{r-1}):0\le p_i\le s_i,
                              \ p_i\ne p_j\ (i\ne j)\}.
\end{equation}
Selecting positions in nondecreasing weight order proves
\begin{equation}\label{eq:statecount}
 |\states(s)|=\prod_{i=0}^{r-1}(s_i+1-i)=\prod_{i=0}^{r-1}q_i.
\end{equation}
All earlier chosen positions lie in the next token's allowed interval, so
exactly $i$ choices are excluded at stage $i$.

The dynamics on \eqref{eq:states} is autonomous. If token $h$ has position
zero, its new position is $s_h$, and another position $p_j$ decreases by one
exactly when $p_j\le s_h$. If no selected token is at zero, write the leader
symbol as $\ast$. Its weight is at least every selected weight, so all
selected positions decrease by one. In a closed core there is no $\ast$.

\begin{lemma}[Explicit inverse]\label{lem:inverse}
The bounded-state map is a permutation of $\states(s)$. Given a state, find
the selected tokens satisfying $p_i=s_i$. If none exists, its predecessor
is $(p_i+1)_i$. Otherwise choose among them a token $h$ with least weight,
put it at zero, increase by one every other position strictly less than
$s_h$, and leave the larger positions unchanged.
\end{lemma}
\begin{proof}
If no equality $p_i=s_i$ holds, every position may be increased by one
without exceeding its bound; this is the inverse of a $\ast$ move.
No selected throw could have produced the state, since that would leave
the thrown token at its weight.

If $h$ has the least weight among matching tokens, increasing a position
$p_j<s_h$ cannot violate $p_j\le s_j$. A violation would require
$p_j=s_j<s_h$, giving a matching token of smaller weight. Thus the proposed
predecessor is valid, and a throw of $h$ recovers the given state.

No other predecessor is possible. A $\ast$ predecessor would exceed a
matching token's bound. A throw from a larger matching weight would likewise
push the smallest matching token beyond its bound when reversed. Matching
tokens cannot tie in weight, since they would occupy the same position.
This proves existence and uniqueness of the predecessor. In a closed core,
a matching token always exists: the occupant of position $r-1$ must have
weight $r-1$.
\end{proof}

\subsection{Deleting a minimum-weight clock}
\begin{lemma}[Minimum deletion]\label{lem:deletion}
In a bounded system, a selected token of minimum weight $m$ leads once every
$q=m+1$ steps. After deleting that token and its leader times, the survivor
process is the same bounded throwback process with all survivor weights
reduced by one.
\end{lemma}
\begin{proof}
Between its throws the minimum token has a position in $\{1,\ldots,m\}$.
Every other leader has weight at least $m$, so each intervening move shifts
it left exactly one place. It therefore has constant return time $m+1$.

When another token of weight $v$ leads, the minimum token lies between
positions $1$ and $m$, and hence is crossed by that throw. Removing it
reduces the throw distance from $v$ to $v-1$. A move of the minimum token
changes no relative order among survivors and disappears in compressed time.
At an instant when the minimum is the head, deleting it reduces every
survivor position by one, placing the remaining selected tokens within their
new bounds. The same reasoning applies to $\ast$ moves.
\end{proof}
The closed recursion ends with a single token of reduced weight zero. Its
word is constant and has period one. This is only a terminal bookkeeping
system; the original queue still has strictly positive weights.

\begin{theorem}[Exact periods, frequencies, and cycle counts]\label{thm:cycles}
For either bounded system above, define
\begin{equation}\label{eq:frequencies}
 R_0=1,\qquad R_{i+1}=R_i\frac{q_i-1}{q_i},\qquad
 f_i=\frac{R_i}{q_i}.
\end{equation}
Every cycle has the same least labeled state period $P(s)$. It is also the
least period of the projected leader word (with a single symbol $\ast$ for
unselected leaders). It is determined by
\begin{equation}\label{eq:period-recurrence}
 P_r=1,\qquad
 P_i=\frac{q_iP_{i+1}}{\gcd(P_{i+1},q_i-1)}
       \quad(i=r-1,r-2,\ldots,0),\qquad P(s)=P_0.
\end{equation}
The frequency of token $i$ on every cycle is $f_i$, and the frequency of
$\ast$ is $R_r$ in the open case. Moreover,
\begin{equation}\label{eq:denominatorperiod}
 P(s)=\lcm_{0\le i<r}\den(f_i),\qquad
 \#\{\text{cycles}\}=\frac{\prod_iq_i}{P(s)}.
\end{equation}
The denominator is taken after reducing each rational number.
\end{theorem}

\begin{proof}
Lemma~\ref{lem:inverse} makes all bounded orbits periodic from the outset.
Delete the minimum token using Lemma~\ref{lem:deletion}. By induction, suppose
the reduced projected word has least period $P'$. Reinsert the minimum
symbol at its fixed residue class modulo $q=q_0$. A period $H$ of the full
word must satisfy $q\mid H$, since the occurrence set of this unique labeled
symbol is a single residue class modulo $q$. Write $H=q\ell$. The same
shift advances the reduced word by $(q-1)\ell$ symbols. Hence it is a period
exactly when
\[
 P'\mid(q-1)\ell.
\]
The least such $H$ is $qP'/\gcd(P',q-1)$. The empty open projection has the
constant word $\ast$ and period one. The single zero-weight closed terminal
system also has period one. Induction proves \eqref{eq:period-recurrence}.

The minimum symbol occupies a proportion $1/q$ of the times. All survivor
frequencies are multiplied by $(q-1)/q$. Iterating gives
\eqref{eq:frequencies}, including the residual $\ast$ frequency.

To verify that the leader word and the state have the same least period,
extend the finite permutation orbit to all integer times. Every selected
token occurs in the word, as its frequency above is positive. To determine
its position at a time $t$, start just after any earlier occurrence of that
token, where its position equals its weight, and replay the projected word.
A $\ast$ move has a completely specified effect on its position. A period
of the projected word thus gives the same position for every selected token
after the shift. The state period divides the word period; the converse is
immediate.

Finally, prove the denominator formula by the same induction. For an integer
$H$ to make all $Hf_i$ integral, the first frequency $1/q$ requires $H=q\ell$.
The survivor frequencies then require that $(q-1)\ell$ times every reduced
frequency be integral, equivalently $P'\mid(q-1)\ell$ by induction. The
residual frequency imposes no extra condition because it is one minus the
sum of the selected frequencies. Thus the least denominator-clearing $H$
is exactly the period already found. All states lie on equal-length cycles,
so their number is the state count divided by that length.
\end{proof}

\section{A complete coprimality test for perfect mixing}\label{sec:mixing}
Perfect mixing of a finite selected family means that its orbit visits every
bounded configuration before returning. Theorem~\ref{thm:cycles} turns this
into an explicit arithmetic test.

\begin{corollary}[Perfect-mixing criterion]\label{cor:mixing}
The bounded system with strides $q_0,\ldots,q_{r-1}$ is a single cycle if
and only if
\begin{equation}\label{eq:mixing}
 \boxed{\quad\gcd(q_i-1,q_j)=1
                  \quad\text{for every }0\le i<j<r.\quad}
\end{equation}
For an increasing infinite input with weights $s_i$, perfect mixing of every
finite initial selected family is therefore equivalent to
\[
 \gcd(s_i-i,\ s_j-j+1)=1\qquad(i<j).
\]
\end{corollary}
\begin{proof}
The period recursion has the form $P_0=\prod_iq_i/\prod_id_i$, where each
$d_i$ is the positive gcd appearing at its stage. Equality with the full
state count requires every $d_i=1$. Working backward, this is equivalent to
\[
 \gcd\Bigl(q_i-1,\prod_{j>i}q_j\Bigr)=1
 \quad\text{for every }i,
\]
which is exactly \eqref{eq:mixing}. For an increasing input, the initial
positions $i$ already satisfy the bounds, so there is no entry transient in
the selected configuration system.
\end{proof}
For consecutive weights beginning with $a$, all strides equal $a+1$ and
all gcds are $\gcd(a,a+1)=1$. This recovers the consecutive-input case of
\cite[Theorem~9]{DD}.

\begin{center}
\begin{tabular}{@{}llllr@{}}
\toprule
Selected weights & Strides $q_i$ & State count & Period & Cycles\\
\midrule
$(1,3,4)$ & $(2,3,3)$ & $18$ & $18$ & $1$\\
$(2,3,5)$ & $(3,3,4)$ & $36$ & $9$ & $4$\\
$(2,4)$ & $(3,4)$ & $12$ & $6$ & $2$\\
$(3,6)$ & $(4,6)$ & $24$ & $8$ & $3$\\
$(2,2)$ & $(3,2)$ & $6$ & $3$ & $2$\\
\bottomrule
\end{tabular}
\end{center}
For example, $(2,3,5)$ fails because
$\gcd(q_0-1,q_2)=\gcd(2,4)=2$. Its four cycles all have length nine.
The criterion makes no claim that merely having pairwise coprime weights is
sufficient; the relevant quantities are the shifted strides in
\eqref{eq:mixing}.

\section{Infinite frequencies, repetitions, and escaped density}\label{sec:densities}
Assume the infinite input is admissible. It has a well-defined sorted labeled
list $s_0\le s_1\le\cdots$, breaking equal-weight ties by original label.
Every finite initial family eventually enters its bounded configuration
space, because all its tokens eventually lead. Theorem~\ref{thm:cycles}
therefore proves, without exchanging an infinite sum and a limit, that the
limiting frequency of sorted token $i$ is
\begin{equation}\label{eq:infinitefreq}
 f_i=\frac1{s_i-i+1}
       \prod_{j<i}\frac{s_j-j}{s_j-j+1}.
\end{equation}
For distinct weights this recovers the formula in \cite[Theorem~4]{DD};
minimum deletion also justifies the repeated-weight case.

\subsection{Equal weights have equal individual frequencies}
Suppose a weight $v$ occurs at sorted indices $a,\ldots,a+b-1$. Then
$q_{a+t}=v-a-t+1$. Telescoping gives
\[
 R_{a+t}=R_a\frac{v-a-t+1}{v-a+1},\qquad
 f_{a+t}=\frac{R_a}{v-a+1}\quad(0\le t<b).
\]
Thus tie-breaking does not affect individual frequencies. The numerical
frequency of the value $v$ is $bR_a/(v-a+1)$.
For $(2,2,3,4,\ldots)$ the individual frequencies begin
\[
 \frac13,\frac13,\frac16,\frac1{12},\frac1{24},\ldots.
\]

\subsection{When do these frequencies sum to one?}
Since $f_i=R_i-R_{i+1}$, finite telescoping gives
\begin{equation}\label{eq:mass}
 \sum_{i=0}^{n-1}f_i=1-R_n,\qquad
 \sum_{i\ge0}f_i=1-R_\infty,\qquad
 R_\infty=\prod_{i\ge0}\left(1-\frac1{q_i}\right).
\end{equation}
For $0\le u\le1/2$,
$u\le-\log(1-u)\le2u$. Since $q_i\ge2$, this proves
\begin{equation}\label{eq:escapecriterion}
 R_\infty=0\quad\Longleftrightarrow\quad
       \sum_{i\ge0}\frac1{q_i}=\infty.
\end{equation}
Existence of every individual frequency alone does not imply that the sum
is one. In particular, the general normalization sentence before Section 3
of \cite[p.~3]{DD} needs this qualification; no infinite sum--limit interchange
is justified merely by pointwise existence.

Here is an exact throwback example. Take
\begin{equation}\label{eq:escapeexample}
 s_i=i+2^{i+1}-1\qquad(i\ge0),
\end{equation}
so $q_i=2^{i+1}$. The input is strictly increasing and admissible, yet
\begin{align*}
 R_\infty&=\prod_{k\ge1}(1-2^{-k})
       =0.288788095086602421278899721929\ldots,\\
 \sum_i f_i&=1-R_\infty
       =0.711211904913397578721100278070\ldots.
\end{align*}
Every token has positive frequency and recurs, but a positive amount of
density escapes to unbounded token ranks. There is no missing leader at any
finite time; rather, density is not continuous under a decreasing countable
intersection of these tail sets.

The displayed digits can be certified with rational arithmetic. If
$R_N=\prod_{k=1}^N(1-2^{-k})$, the elementary inequality
$\prod_k(1-u_k)\ge1-\sum_ku_k$ gives
\begin{equation}\label{eq:productbound}
 R_N(1-2^{-N})\le R_\infty\le R_N.
\end{equation}
The companion data use $N=160$, making the interval width less than $2^{-160}$.

\subsection{Uniform recurrence for nondecreasing inputs}
\begin{proposition}\label{prop:uniform}
For every nondecreasing positive input, including repeated weights, every
finite block of its labeled leader word reappears with bounded gaps.
\end{proposition}
\begin{proof}
In the admissible case, the first $r$ tokens start at positions $i\le s_i$,
so their bounded projected word is periodic from time zero. Any finite block
of the full word uses only finitely many labels; choose $r$ containing them
all. At that block, the projected word has no $\ast$ symbols. Repetition by
its period reproduces the full block, with a fixed gap bound.

In the obstructed case, a nondecreasing input's first bad prefix is itself
critical. Indeed, if its first failure occurs at index $J$, then
$x_{J-1}\ge J$ when $J\ge1$, whereas failure forces $x_J\le J$; hence
$x_J=J$ and all $J+1$ initial weights are at most $J$. This prefix is closed
and all its initial positions satisfy their weight bounds. Its labeled word
is therefore periodic from time zero by Theorem~\ref{thm:cycles}.
\end{proof}
For unsorted admissible inputs, finite small-weight projections are eventually
periodic, but a finite transient block need not recur. The proposition is
not silently asserting uniform recurrence for every arbitrary ordering.

\section{The largest possible recurrent-core period}\label{sec:extremal}
Let $M_r$ be the largest \emph{labeled} eventual period among critical cores
of size $r\ge2$. Although a core can have up to $r!$ possible labeled orders,
its dynamics has a much sharper extremal period.

\begin{theorem}[Sharp extremal period]\label{thm:extremal}
For every $r\ge2$,
\begin{equation}\label{eq:maxperiod}
 \boxed{\qquad M_r=\max_{2\le K\le r}K^{r-K+1}.\qquad}
\end{equation}
For any candidate $K$, the critical profile
\begin{equation}\label{eq:extremalprofile}
 s_i=\min(K+i-1,r-1),\qquad 0\le i<r,
\end{equation}
has period exactly $K^{r-K+1}$.
\end{theorem}
\begin{proof}
Reverse the stride list of a critical profile, writing
$a_t=q_{r-1-t}$ for $0\le t<r$. Criticality and sortedness imply
\[
 a_0=1,\qquad a_t\ge2\ (t\ge1),\qquad a_{t+1}\le a_t+1.
\]
Let $p_t$ be the period accumulated through $a_t$ in the backward recursion.
Then $p_0=1$, each $p_t$ is a multiple of $a_t$, and
\[
 p_{t+1}=\frac{a_{t+1}p_t}{\gcd(p_t,a_{t+1}-1)}.
\]
A unit rise from $a$ to $a+1$ is special: since $a\mid p_t$,
\begin{equation}\label{eq:rise}
 p_{t+1}=\frac{a+1}{a}p_t.
\end{equation}
Thus consecutive unit rises telescope instead of multiplying their heights.

Partition the path $a_0,\ldots,a_{r-1}$ into maximal runs of consecutive
unit rises. The first run, ending at height $h$, changes the period from
$1$ to $h$. At the first step of any later run, from a height $a$ to
$b\le a$, the period grows by at most the factor $b$. Subsequent unit rises
in that run telescope, so the entire run grows the period by at most its
final height $h$.

Let $K=\max_ta_t$, let $A$ be the number of unit-rise edges, and let $B$
be the number of runs. Then $B=r-A$, while reaching $K$ from $1$ requires
$A\ge K-1$. Every run ends at height at most $K$. Therefore
\[
 P(s)\le K^B\le K^{r-K+1}.
\]

For equality at a prescribed $K$, use the reverse stride path
\[
 (1,2,\ldots,K,\underbrace{K,\ldots,K}_{r-K\text{ further terms}}).
\]
The first run has period $K$. Every later constant-$K$ step multiplies the
period by $K$, since a power of $K$ is coprime to $K-1$. Its period is
$K^{r-K+1}$. Reversing the strides gives precisely
\eqref{eq:extremalprofile}, which is critical. Taking the maximum proves
\eqref{eq:maxperiod}.
\end{proof}

The values for $r=2,3,\ldots$ begin
\[
 2,4,9,27,81,256,1024,4096,16384,78125,390625,1953125,\ldots.
\]
OEIS A003320 is defined by $a(n)=\max_{0\le k\le n}k^{n-k}$
\cite{A003320}. Thus
\begin{equation}\label{eq:oeismax}
 M_r=a(r+1)\qquad(r\ge2).
\end{equation}
The extremal formula itself is a known integer sequence; its interpretation
as the largest throwback core period is the result established here.

\subsection{An exact two-candidate rule and a Lambert-\texorpdfstring{$W$}{W} asymptotic}
Write $N=r+1$ and maximize
\[
 F_N(x)=(N-x)\log x.
\]
Since
\[
 F_N'(x)=\frac Nx-1-\log x,
 \qquad F_N''(x)=-\frac N{x^2}-\frac1x<0,
\]
there is a unique continuous maximum. Let $w=W(eN)$ be the positive real
solution of $we^w=eN$. The maximizing point is
\begin{equation}\label{eq:kstar}
 x_*=\frac Nw=e^{w-1}.
\end{equation}
For $r\ge3$, a maximizing integer is one of
$\lfloor x_*\rfloor$ and $\lceil x_*\rceil$, so comparing those two exact
integer powers evaluates $M_r$. The case $r=2$ is $M_2=2$.

At the continuous optimum,
\begin{equation}\label{eq:continuous-max}
 F_N(x_*)=N\left(w-2+\frac1w\right).
\end{equation}
For sufficiently large $N$, rounding to the nearest integer changes $x$ by
at most $1/2$. Taylor's theorem and the displayed second derivative give
\begin{equation}\label{eq:rounding}
 0\le N\left(w-2+\frac1w\right)-\log M_r
       =O\left(\frac{w^2}{N}\right).
\end{equation}
In particular, the error here tends to zero, so this is also a relative
asymptotic for $M_r$ after exponentiating:
\begin{equation}\label{eq:relative-asymptotic}
 M_r=\exp\!\left[N\left(w-2+\frac1w\right)\right]
       \left(1+O\left(\frac{(\log N)^2}{N}\right)\right).
\end{equation}
The rounded value never exceeds the continuous one; the unsigned $O$ notation
in \eqref{eq:relative-asymptotic} does not discard that fact.

To make the scale more explicit, put $L=\log N$ and $\ell=\log L$.
Solving $w+\log w=L+1$ by substitution yields
\[
 w=L-\ell+1+\frac{\ell-1}{L}
       +\frac{\ell^2-4\ell+3}{2L^2}
       +O\left(\frac{\ell^3}{L^3}\right).
\]
Substituting in \eqref{eq:rounding} gives
\begin{equation}\label{eq:logexpansion}
 \log M_r=N\left[L-\ell-1+\frac\ell L
                 +\frac{(\ell-1)^2}{2L^2}
                 +O\left(\frac{\ell^3}{L^3}\right)\right]
            +O\left(\frac{L^2}{N}\right).
\end{equation}
The remainder for $w$ follows rigorously because substituting its displayed
truncation into $w+\log w$ leaves a remainder of the stated order, and the
derivative $1+1/w$ is bounded below by one. Formula~\eqref{eq:logexpansion}
is a logarithmic expansion: dropping its subleading terms and claiming a
relative asymptotic for the unlogged sequence would not be valid.

\section{Rational generating functions and exact simulation}\label{sec:implementation}
Once a recurrent core is known, every numerical or labeled leader generating
function is rational. More generally, a finite selected projection in the
admissible regime is eventually periodic. If its entry time is $t_0$, its
period is $P$, and token $i$ occurs at the residue set
$B_i\subseteq\{0,\ldots,P-1\}$ on that periodic part, then
\begin{equation}\label{eq:leadergf}
 \sum_{t\ge0}\ind_{\{L_t=i\}}z^t
   =\sum_{t<t_0}\ind_{\{L_t=i\}}z^t
     +\frac{z^{t_0}\sum_{b\in B_i}z^b}{1-z^P}.
\end{equation}
The same formula, with coefficients equal to leader weights, gives the
numerical leader series of an obstructed input. It need not hold for the
entire infinite-alphabet word in the admissible alternative.

For \eqref{eq:example100}, the numerical word has the generating function
\begin{equation}\label{eq:examplegf}
 2+100z+2z^2+2z^3+\frac{z^4+2z^5}{1-z^2}.
\end{equation}
The period-two tail here does not contradict the labeled period four.

\subsection{A finite exact model from a prefix certificate}
There is no need to materialize an infinite tail, or even a queue of length
the largest throw. Given the obstruction certificate, retain only the
positions of labels $0,\ldots,J$, initially $(0,1,\ldots,J)$. A token from
this set is always at zero, because Theorem~\ref{thm:obstruction} says no
later label ever leads. A move changes these finitely many positions using
its known weight. Each coordinate is bounded by $\max(i,x_i)$, so this is a
finite exact state model of the \emph{whole} leader word.

The function \texttt{certified\_orbit} implements this model and detects its
first repeated state. Its optional step cap is a resource limit; reaching
it does not invalidate the static certificate or constitute a recurrence
counterexample. Static core extraction and the formula for $P$ can be fast
even when explicit enumeration of a period is necessarily huge.

\subsection{Avoiding a truncation error}
A finite prefix is not automatically a closed queue. In particular, Python's
\texttt{list.insert(v, head)} silently appends if $v$ exceeds the list length;
using that operation on an insufficient prefix changes the process.

The function \texttt{finite\_horizon} instead starts with at least $T+w$
input tokens to compute $T$ moves and a final visible window of length $w$.
A throw beyond the represented prefix is omitted rather than incorrectly
appended. Such a token, and every unrepresented tail token, can move left by
at most one per remaining step, so cannot reach the requested window in time.
The representation loses at most one position per move and remains long
enough for the promised horizon. This is an exact finite-horizon computation,
not an empirical tail approximation.

\section{Computational verification and reproducibility}\label{sec:verification}
All distributed Python code uses only the standard library. The following
checks were actually executed; the raw counts, example words, and assertions
are included in \texttt{data/}.

\subsection{Every small closed input, with repeated weights distinguished}
For each $2\le n\le7$, every ordered word in $\{1,\ldots,n-1\}^n$ was
simulated as a genuinely closed labeled queue until a state repeated. The
check compared the recurrent labels and the labels ever visited with the
static first-obstruction certificate, verified each individual rational
frequency, checked the exact period, and confirmed that the tail was fixed
on the eventual cycle.

\begin{center}
\begin{tabular}{@{}rrrrrr@{}}
\toprule
$n$ & Words checked & Critical words & Critical multisets & Max. period & Max. transient\\
\midrule
2&1&1&1&2&0\\
3&8&4&2&4&2\\
4&81&27&5&9&6\\
5&1024&256&14&27&14\\
6&15625&3125&42&81&30\\
7&279936&46656&132&256&62\\
\bottomrule
\end{tabular}
\end{center}
The total is $296\,675$ closed inputs. The transient column is an observed
finite-range statistic, not a separately asserted sharp transient formula.
The maximum-period column agrees with Theorem~\ref{thm:extremal}.

\subsection{Complete bounded configuration spaces}
For every admissible nondecreasing selected profile with weights at most
seven and at most five tokens, all states in \eqref{eq:states} were enumerated.
The test verified the explicit inverse at every state, decomposed the entire
permutation into cycles, and checked both periods and frequencies on each
cycle.

\begin{center}
\begin{tabular}{@{}rrrrrr@{}}
\toprule
Selected tokens & Profiles & States & Cycles & One-cycle profiles & Max. period\\
\midrule
1&7&35&7&7&8\\
2&27&600&53&17&49\\
3&75&6500&358&19&216\\
4&165&48181&1865&25&625\\
5&297&247929&8025&11&1024\\
\bottomrule
\end{tabular}
\end{center}
This covers $571$ profiles, $303\,245$ bounded states, and $10\,308$ cycles.
It tests the coprimality criterion against actual cycle decompositions, rather
than merely comparing two algebraically equivalent formula implementations.

\subsection{Additional independent checks}
The waiting-or-barrier bound was checked in $5470$ finite token-position
cases. Sparse backward parking was compared with a sorted-prefix test on
$1500$ seeded random inputs, some containing weights around $10^{40}$.
The finite-horizon simulator was compared with sufficiently long closed
reference queues on $200$ cases. Separately, every critical stride profile
through core size $13$ was enumerated---$290\,511$ profiles in total---and its
period was compared with the sharp extremal formula. Finally, the escape
product was bounded using exact fractions as in \eqref{eq:productbound}.

These are consistency checks over explicitly finite ranges. They are not
Lean verification, exhaustive checks of all infinite inputs, or independent
human review of the proofs.

\subsection{Reproduction commands}
From the archive's top directory, run
\begin{lstlisting}
python3 code/verify.py
python3 code/extra_results.py
pdflatex -interaction=nonstopmode -halt-on-error throwback_research.tex
pdflatex -interaction=nonstopmode -halt-on-error throwback_research.tex
\end{lstlisting}
The Python scripts write deterministic mathematical outputs, except for the
recorded wall-clock runtime. Exact frequencies are stored as rational numbers
or rational strings, not inferred from floating-point histograms.

\section{Scope, remaining refinements, and conclusion}\label{sec:scope}
The repeated-entry question has a complete answer at the labeled level:
a simple global counting condition decides universal recurrence, and the
first bad prefix specifies the exact finite survivor set when it fails.
Corollary~\ref{cor:numerical} translates this to numerical values. The
bounded-state analysis also yields a necessary and sufficient perfect-mixing
criterion, rather than a collection of sufficient examples.

Several refinements are distinct from these resolved statements. With equal
weights, collapsing token labels can reduce the period in ways that depend
on the cycle. For example, selected weights $(3,3)$ have labeled projected
period four, but their numerical projection can be $3,3,\ast,\ast$ of period
four or $3,\ast,3,\ast$ of period two. Thus numerical period is not in general
determined by the common numerical frequencies alone. The article does not
claim a complete enumeration of these quotient cycles. Nor does it claim a
sharp formula for every capture transient, or a bit-optimal implementation
for very large integer weights.

The recurrence and cycle proofs are elementary enough to suggest a compact
formalization: finite positions, nested count inequalities, a finite
permutation, and gcd arithmetic carry the main argument. No such
formalization is included or claimed. The appropriate next validation step
is independent checking of the proofs and a broader priority search before
any submission or OEIS amendment.

The resulting picture is unusually explicit. A queue either keeps every
occurrence active, or a finite prefix certificate predicts its permanent
core. Finite projections consist of equal-length cycles controlled by simple
clocks. Their exact periods identify when perfect mixing occurs, and the
largest core period is a classical maximum-of-powers sequence with a precise
Lambert-$W$ scale. These deductions connect recurrence dynamics, parking
functions, Catalan counting, rational generating functions, and asymptotic
analysis in a single framework.

\appendix
\section{A compact proof-audit checklist}\label{app:audit}
The following points are the main places where an incorrect generalization
could otherwise enter.

\paragraph{Occurrences are never identified.}
All positional arguments concern distinct labels, even when their weights
coincide. Only after the recurrent labels are known are frequencies grouped
by numerical value.

\paragraph{The obstruction is an initial-prefix obstruction.}
The smallest overloaded threshold of the whole input need not identify the
actual core. Both the least prefix index $J$ and then the least threshold
inside that prefix are required.

\paragraph{Capture is proved in finite time.}
The argument does not assume that all globally small tokens eventually
surface. Only tokens in the first bad prefix are forced to surface, using
the admissibility of every earlier prefix and the waiting-or-barrier lemma.

\paragraph{The period is a least period.}
Periodicity alone would provide only an upper bound. The unique minimum
symbol forces divisibility by its clock length, and the compressed survivor
word forces the remaining divisibility condition. This proves minimality.

\paragraph{A frequency denominator alone need not give a period.}
The lcm formula is established by the clock recursion, not by the generally
false inference that frequencies determine every periodic word. Numerical
quotients explicitly demonstrate that distinction.

\paragraph{Infinite density is handled by finite telescoping first.}
Equation~\eqref{eq:mass} is obtained from finite sums before passing to a
monotone product limit. No unjustified exchange of a countable sum and a
time-average limit is used.

\paragraph{Enumeration and numerical experiments have separate roles.}
The Catalan and power counts are proved by cyclic minima. Enumerations check
those proofs on finite ranges, and do not stand in for an all-size argument.

\section{Code and data inventory}\label{app:files}
\begin{longtable}{@{}p{0.35\textwidth}p{0.60\textwidth}@{}}
\toprule
File & Purpose\\
\midrule
\endhead
\texttt{code/throwback.py} & Sparse static classifier, independent certificate
verifier, bounded map and inverse, exact rational statistics, safe finite
horizon, and certified prefix orbit.\\
\texttt{code/verify.py} & Exhaustive closed queues and bounded configuration
spaces, plus waiting, sparse, and horizon checks.\\
\texttt{code/extra\_results.py} & Exhaustive critical profiles through size
13, worked examples, and rational escape-product bounds.\\
\texttt{code/analyze.py} & Command-line interface for prefix certificates
and exact orbit analysis.\\
\texttt{data/verification.json} & Executed verification ranges and exact totals.\\
\texttt{data/finite\_queue\_checks.csv} & Closed-input enumeration table.\\
\path{data/bounded_configuration_checks.csv} & Full finite permutation checks.\\
\texttt{data/extremal\_periods.csv} & Maximum periods and attaining profiles.\\
\texttt{data/examples.json} & Exact transient and periodic example words.\\
\path{data/escape_product_bounds.json} & Certified rational-product interval,
displayed with high-precision decimal endpoints.\\
\texttt{notes/source\_status.md} & Sources consulted and the scope of the
novelty/status check.\\
\bottomrule
\end{longtable}

\begin{thebibliography}{9}
\bibitem{DD}
J. Darnell and B. F. Dribus,
\emph{Throwback Sequences of Positive Integers},
Journal of Integer Sequences \textbf{28} (2025), Article 25.5.1, 18 pp.
\url{https://cs.uwaterloo.ca/journals/JIS/VOL28/Dribus/dribus4.pdf}.

\bibitem{A357081}
OEIS Foundation Inc., \emph{The On-Line Encyclopedia of Integer Sequences},
A357081, ``Leader at step $n$ of the THROWBACK procedure.''
\url{https://oeis.org/A357081}. Accessed 19 September 2026.

\bibitem{A000108}
OEIS Foundation Inc., \emph{The On-Line Encyclopedia of Integer Sequences},
A000108, Catalan numbers.
\url{https://oeis.org/A000108}. Accessed 19 September 2026.

\bibitem{A000312}
OEIS Foundation Inc., \emph{The On-Line Encyclopedia of Integer Sequences},
A000312, $a(n)=n^n$, including prime parking functions of length $n+1$.
\url{https://oeis.org/A000312}. Accessed 19 September 2026.

\bibitem{A003320}
OEIS Foundation Inc., \emph{The On-Line Encyclopedia of Integer Sequences},
A003320, $a(n)=\max_{0\le k\le n}k^{n-k}$.
\url{https://oeis.org/A003320}. Accessed 19 September 2026.
\end{thebibliography}
\end{document}
